%% 06-test-suite.tex — Test suite and Vermont walkthrough

\section{Test Suite and Vermont Walkthrough}
\label{sec:testsuite}

\subsection*{Quick Health Check: \texttt{redist doctor}}

Before running any full-scale replication, confirm the environment is
correctly configured:

\begin{verbatim}
redist doctor                        # General environment check
redist doctor --check-tutorial-data  # Verify Vermont fixture data
\end{verbatim}

\noindent The \texttt{doctor} command checks:
\begin{itemize}[itemsep=2pt]
  \item Rust version ($\geq$ 1.77)
  \item METIS availability (vendored, always passes)
  \item Python 3.13+ (for dashboard generation)
  \item Census data cache status
  \item Vermont walkthrough fixture integrity
\end{itemize}

\subsection*{Vermont 2020 Canonical Walkthrough}

Vermont is the smallest state with a single congressional district,
making it the fastest test case.
However, the walkthrough fixture is artificially expanded to
\textit{simulate} a multi-district state for testing purposes,
using Vermont's tract data with a synthetic two-seat configuration.

\begin{verbatim}
# Run the walkthrough (< 2 minutes on any modern laptop)
redist state --state VT --year 2020 --version walkthrough

# Verify the output matches the fixture
redist label-verify walkthrough

# Browse the output
ls outputs/walkthrough/2020/vermont/
\end{verbatim}

\noindent Expected outputs for the Vermont walkthrough:

\begin{table}[h]
\centering
\caption{Expected Vermont walkthrough outputs (placeholder hashes).}
\label{tab:vt}
\begin{tabular}{@{}lll@{}}
\toprule
File & Description & SHA-256 (placeholder) \\
\midrule
\texttt{districts.csv}   & Tract-to-district assignments & \texttt{abcdef01...} \\
\texttt{summary.csv}     & Compactness, population stats & \texttt{23456789...} \\
\texttt{map\_round1.svg} & Bisection round 1 map         & \texttt{fedcba98...} \\
\texttt{map\_final.svg}  & Final district map            & \texttt{76543210...} \\
\bottomrule
\end{tabular}
\end{table}

\noindent These hashes are pinned by \texttt{examples/vermont-2020-walkthrough/pin.sh}
after the first clean run~\cite{vermont2020walkthrough}.

\subsection*{Full Test Suite}

The repository includes over 1,100 automated tests:

\begin{verbatim}
cargo test --workspace          # All Rust unit and integration tests
pytest tests/unit/ -v           # Python unit tests
pytest tests/integration/ -v    # Integration tests (require pipeline outputs)
\end{verbatim}

\noindent Expected test counts and runtimes:

\begin{table}[h]
\centering
\caption{Test suite summary.}
\label{tab:tests}
\begin{tabular}{@{}llrr@{}}
\toprule
Suite & Language & Tests & Typical runtime \\
\midrule
Rust unit tests       & Rust   & $\sim$900 & $<$30 s \\
Python unit tests     & Python & $\sim$200 & $<$60 s \\
Integration tests     & Python & $\sim$50  & 2--5 min (requires data) \\
\midrule
\textbf{Total}        &        & $\sim$1150 & $<$8 min \\
\bottomrule
\end{tabular}
\end{table}

\subsection*{Expected Runtimes for Full Replication}

\begin{table}[h]
\centering
\caption{Expected runtimes on reference hardware (AMD Ryzen 9, 32 GB RAM).}
\label{tab:runtimes}
\begin{tabular}{@{}lll@{}}
\toprule
Task & States & Time \\
\midrule
Vermont only (test)        & 1  & 30 s -- 2 min \\
Single large state (CA)    & 1  & 3--6 min \\
All 50 states, 2020        & 50 & $\sim$18 min \\
All 50 states, 3 years     & 150 & 2--4 h \\
\bottomrule
\end{tabular}
\end{table}
